\section{Projects for Chapter 0: Introduction}

The introductory chapter is about modelling discipline. The most useful project is therefore an interactive map of the objects that appear between a real situation and a probability calculation.

\subsection*{Project: model-construction map}

The application should use several short scenarios, such as coin tosses, dice, cards, an urn, a waiting time, and a continuous measurement. For each scenario, the student should classify descriptions as
\[
\begin{aligned}
&\text{experiment},\quad \text{recording rule},\quad \text{outcome},\\
&\text{sample space},\quad \text{event},\quad \text{probability law},\\
&\text{random variable},\quad \text{data},\quad \text{statistic}.
\end{aligned}
\]

\begin{examplebox}
\textbf{Prompt to give an AI coding assistant}

\begin{quote}\small
Build an interactive probability-model map. Include scenarios based on coins, dice, cards, urns, waiting times, and a continuous interval. Present cards describing the experiment, recording rule, elementary outcome, sample space, event, probability assignment, random variable, observed data, statistic, and possible conclusion.

Let the student arrange the cards in a logical order. Check the arrangement and explain every incorrect placement. Emphasize the difference between a theoretical probability law, one observed dataset, and a sampling distribution. Add controls that reveal one correct construction and load another scenario.
\end{quote}
\end{examplebox}

A useful extension is an assumption panel containing independence, equal likelihood, replacement, identical distribution, finite variance, and representative sampling. When an assumption is removed, the application should state which later conclusions may fail.

\begin{summarybox}
The project should make the modelling sequence visible:
\[
\text{situation}\longrightarrow\text{model}\longrightarrow\text{calculation}\longrightarrow\text{interpretation}.
\]
\end{summarybox}
